\documentclass[../../main.tex]{subfiles}% 注意这里的文件路径不能用 ./main.tex ,否则用latexmk编译子文件会报错
\graphicspath{{\subfix{./image/}}{\subfix{../../image/}}} % 指定图片引用路径,后续可以直接使用图片文件名
% 如果图片存放在上级目录,则使用 ../../image/ 作为备用路径

% 例如:
% \begin{figure}[H]
% \centering
% \includegraphics[scale=0.3]{图.png}
% \caption{}
% \label{figure:图}
% \end{figure}
% 注意:上述\label{}一定要放在\caption{}之后,否则引用图片序号会只会显示??.

\begin{document}

\section{常曲率曲面}

\begin{definition}
Gauss曲率为常数的曲面称为\textbf{常曲率曲面}。
\end{definition}

\begin{proposition}
设常曲率曲面在测地平行坐标系$(u,v)$下的第一基本形式为
\begin{align}
\mathrm{I}=(\mathrm{d}u)^2+G(u,v)(\mathrm{d}v)^2,\label{eq:4.1}
\end{align}
且$G(u,v)$满足
\begin{align}
G(0,v)=1,\quad \frac{\partial G}{\partial u}(0,v)=0,\label{eq:4.2}
\end{align}
则$\sqrt{G}$满足二阶常系数线性齐次微分方程
\begin{align}
(\sqrt{G})_{uu}+K\sqrt{G}=0.\label{eq:4.3}
\end{align}
\end{proposition}

\begin{proof}

由Gauss曲率的内蕴表达式
\begin{align*}
K&=-\frac{1}{\sqrt{EG}}\left(\left(\frac{(\sqrt{E})_v}{\sqrt{G}}\right)_v+\left(\frac{(\sqrt{G})_u}{\sqrt{E}}\right)_u\right),
\end{align*}
代入$E=1,F=0$，得
\begin{align*}
K=-\frac{1}{\sqrt{G}}(\sqrt{G})_{uu},
\end{align*}
整理即得\eqref{eq:4.3}.

\end{proof}

\begin{theorem}
设曲面$S$的Gauss曲率为常数$K$，在测地平行坐标系$(u,v)$下，$S$的第一基本形式为
\begin{enumerate}[(1)]
\item 当$K>0$时，$\mathrm{I}=(\mathrm{d}u)^2+\cos^2(\sqrt{K}u)(\mathrm{d}v)^2$；
\item 当$K=0$时，$\mathrm{I}=(\mathrm{d}u)^2+(\mathrm{d}v)^2$；
\item 当$K<0$时，$\mathrm{I}=(\mathrm{d}u)^2+\cosh^2(\sqrt{-K}u)(\mathrm{d}v)^2$.
\end{enumerate}
\end{theorem}

\begin{proof}

微分方程\eqref{eq:4.3}的特征方程为$\lambda^2+K=0$，结合初始条件
\begin{align}
\sqrt{G}(0,v)=1,\quad (\sqrt{G})_u(0,v)=0,\label{eq:4.4}
\end{align}
分三类求解：
\begin{enumerate}[(1)]
\item 若$K>0$，特征根为$\lambda=\pm\mathrm{i}\sqrt{K}$，通解为$\sqrt{G}=a(v)\cos(\sqrt{K}u)+b(v)\sin(\sqrt{K}u)$，代入初始条件得$a(v)=1,b(v)=0$，故$\sqrt{G}=\cos(\sqrt{K}u)$；
\item 若$K=0$，通解为$\sqrt{G}=a(v)+b(v)u$，代入初始条件得$a(v)=1,b(v)=0$，故$\sqrt{G}=1$；
\item 若$K<0$，特征根为$\lambda=\pm\sqrt{-K}$，通解为$\sqrt{G}=a(v)\cosh(\sqrt{-K}u)+b(v)\sinh(\sqrt{-K}u)$，代入初始条件得$a(v)=1,b(v)=0$，故$\sqrt{G}=\cosh(\sqrt{-K}u)$.
\end{enumerate}
代入\eqref{eq:4.1}即证.

\end{proof}

\begin{theorem}
具有相同常数Gauss曲率$K$的任意两块常曲率曲面在局部上必可建立保长对应。
\end{theorem}

\begin{proof}

相同$K$的常曲率曲面在测地平行坐标系下第一基本形式完全一致，故局部保长。

\end{proof}

\begin{definition}
设$D$为$uv$平面上的区域，在$D$上指定正定二次微分形式
\begin{align}
\mathrm{d}s^2=E(u,v)(\mathrm{d}u)^2+2F(u,v)\mathrm{d}u\mathrm{d}v+G(u,v)(\mathrm{d}v)^2,\label{eq:4.5}
\end{align}
则称二元对$(D,\mathrm{d}s^2)$为\textbf{抽象曲面}，$\mathrm{d}s^2$称为该抽象曲面的\textbf{度量形式}。
\end{definition}

\begin{proposition}
对抽象曲面$(D,\mathrm{d}s^2)$，有：
\begin{enumerate}[(1)]
\item 切向量$(\mathrm{d}u,\mathrm{d}v)$与$(\delta u,\delta v)$夹角$\theta$的余弦为
\begin{align}
\cos\theta=\frac{E(u,v)\mathrm{d}u\delta u+F(u,v)(\mathrm{d}u\delta v+\delta u\mathrm{d}v)+G(u,v)\mathrm{d}v\delta v}{\sqrt{E(\mathrm{d}u)^2+2F\mathrm{d}u\mathrm{d}v+G(\mathrm{d}v)^2}\sqrt{E(\delta u)^2+2F\delta u\delta v+G(\delta v)^2}};\label{eq:4.6}
\end{align}
\item 曲线$u=u(t),v=v(t)\ (a\leqslant t\leqslant b)$的弧长为
\begin{align}
\int_{a}^{b}\sqrt{E\left(\frac{\mathrm{d}u}{\mathrm{d}t}\right)^2+2F\frac{\mathrm{d}u}{\mathrm{d}t}\frac{\mathrm{d}v}{\mathrm{d}t}+G\left(\frac{\mathrm{d}v}{\mathrm{d}t}\right)^2}\mathrm{d}t.\label{eq:4.7}
\end{align}
\end{enumerate}
\end{proposition}

\begin{remark}
1854年，Riemann将Gauss内蕴微分几何推广至任意$n$维，开创\textbf{Riemann几何}，二维Riemann几何即为Gauss内蕴微分几何。抽象曲面的核心几何量为Gauss曲率、测地曲率、测地线。欧氏几何与非欧几何的本质区别为空间弯曲程度：欧氏空间$K=0$，非欧空间$K$为非零常数，弯曲程度决定测地线性状与测地三角形内角和。
\end{remark}

\begin{proposition}
度量形式
\begin{align}
\mathrm{d}s^2=\frac{(\mathrm{d}u)^2+(\mathrm{d}v)^2}{\left(1+\frac{K}{4}(u^2+v^2)\right)^2}\label{eq:4.8}
\end{align}
的Gauss曲率为常数$K$。当$K\geqslant0$时定义域为整个$uv$平面；当$K<0$时定义域为区域
\begin{align*}
D=\left\{(u,v):u^2+v^2<-\frac{4}{K}\right\}.
\end{align*}
\end{proposition}

\begin{definition}
当$K<0$时，区域$D=\left\{(u,v):u^2+v^2<-\frac{4}{K}\right\}$配备度量形式\eqref{eq:4.8}的抽象曲面称为\textbf{Klein圆}。
\end{definition}

\begin{example}
当$K=0$时，$\mathrm{d}s^2=(\mathrm{d}u)^2+(\mathrm{d}v)^2$，该抽象曲面即为普通平面，测地线为普通直线。
\end{example}

\begin{proposition}
当$K>0$时，Gauss曲率为$K$的抽象曲面可视为$\mathbb{E}^3$中半径为$\frac{1}{\sqrt{K}}$的球面经球极投影的像，球极投影变换为
\begin{align}
u=\frac{2x}{\sqrt{K}z+1},\quad v=\frac{2y}{\sqrt{K}z+1},\quad x^2+y^2+z^2=\frac{1}{K},\label{eq:4.9}
\end{align}
逆变换为
\begin{align}
\begin{cases}
x=\frac{4u}{4+K(u^2+v^2)},\\
y=\frac{4v}{4+K(u^2+v^2)},\\
z=\frac{1}{\sqrt{K}}\cdot\frac{4-K(u^2+v^2)}{4+K(u^2+v^2)}.
\end{cases}\label{eq:4.10}
\end{align}
\end{proposition}

\begin{proof}

直接计算可得球面\eqref{eq:4.10}的第一基本形式为\eqref{eq:4.8}。球面上测地线为大圆周，其球极投影为$uv$平面上以原点为中心、半径$\frac{2}{\sqrt{K}}$的圆周，及过该圆周一对对径点的圆周或直线。

\end{proof}

\begin{remark}
$K>0$时的抽象曲面中，任意两条测地线必相交。
\end{remark}

\begin{definition}
\textbf{洛伦兹空间}$L^3$的点为有序三元实数组$(x,y,z)$，向量$\bm{a}=(x_1,y_1,z_1),\bm{b}=(x_2,y_2,z_2)$的内积定义为
\begin{align}
\bm{a}\cdot\bm{b}=x_1x_2+y_1y_2-z_1z_2.\label{eq:4.11}
\end{align}
\end{definition}

\begin{proposition}
洛伦兹空间$L^3$中的曲面
\begin{align}
\Sigma=\left\{(x,y,z)\in L^3:x^2+y^2-z^2=\frac{1}{K},\ z>0\right\}\label{eq:4.12}
\end{align}
经球极投影
\begin{align}
u=\frac{2x}{\sqrt{-K}z+1},\quad v=\frac{2y}{\sqrt{-K}z+1}\label{eq:4.13}
\end{align}
及其逆变换
\begin{align}
\begin{cases}
x=\frac{4u}{4+K(u^2+v^2)},\\
y=\frac{4v}{4+K(u^2+v^2)},\\
z=\frac{1}{\sqrt{-K}}\cdot\frac{4-K(u^2+v^2)}{4+K(u^2+v^2)}.
\end{cases}\label{eq:4.14}
\end{align}
与Klein圆一一对应，且$\Sigma$的诱导第一基本形式为\eqref{eq:4.8}。
\end{proposition}

\begin{proof}

对\eqref{eq:4.14}微分得
\begin{align*}
\mathrm{d}x&=\frac{4(4+K(-u^2+v^2))\mathrm{d}u-8Kuv\mathrm{d}v}{(4+K(u^2+v^2))^2},\\
\mathrm{d}y&=\frac{-8Kuv\mathrm{d}u+4(4+K(u^2-v^2))\mathrm{d}v}{(4+K(u^2+v^2))^2},\\
\mathrm{d}z&=\frac{16\sqrt{-K}(u\mathrm{d}u+v\mathrm{d}v)}{(4+K(u^2+v^2))^2},
\end{align*}
代入洛伦兹内积的诱导度量，得
\begin{align*}
\mathrm{I}=\frac{16((\mathrm{d}u)^2+(\mathrm{d}v)^2)}{(4+K(u^2+v^2))^2}=\frac{(\mathrm{d}u)^2+(\mathrm{d}v)^2}{\left(1+\frac{K}{4}(u^2+v^2)\right)^2},
\end{align*}
即式\eqref{eq:4.8}.

\end{proof}

\begin{theorem}
Klein圆上的测地线在球极投影下为$uv$平面内与区域$D$边界$u^2+v^2=-\frac{4}{K}$正交的圆弧或直径，满足：过“直线”外一点可作无数条“直线”与已知“直线”不相交，即非欧几何平行公理。
\end{theorem}

\begin{proof}

曲面$\Sigma$满足$\bm{r}\cdot\bm{r}=\frac{1}{K}\ (z>0)$，微分得$\mathrm{d}\bm{r}\cdot\bm{r}=0$，即向径$\bm{r}$为$\Sigma$的法向量。
设$L^3$中过原点的平面$\Pi$与$\Sigma$的交线为弧长参数曲线$\bm{r}(s)$，则$\frac{\mathrm{d}\bm{r}}{\mathrm{d}s}\cdot\frac{\mathrm{d}\bm{r}}{\mathrm{d}s}=1$，求导得$\frac{\mathrm{d}^2\bm{r}}{\mathrm{d}s^2}\cdot\frac{\mathrm{d}\bm{r}}{\mathrm{d}s}=0$；结合$\frac{\mathrm{d}\bm{r}}{\mathrm{d}s}\cdot\bm{r}=0$，得$\frac{\mathrm{d}^2\bm{r}}{\mathrm{d}s^2}$与法向量$\bm{r}$平行，故该曲线为$\Sigma$的测地线。
该测地线经球极投影后即为$uv$平面内与$D$边界正交的圆弧或直径，由此可得非欧几何平行公理。

\end{proof}

\begin{example}
试在测地极坐标系下写出常曲率曲面的第一基本形式。
\end{example}

\begin{example}
证明: 在常曲率曲面上, 以任意一点 $p$ 为中心的测地圆的测地曲率为常数。
\end{example}

\begin{example}
已知常曲率曲面的第一基本形式为
\begin{align*}
\mathrm{I}=
\begin{cases}
(\mathrm{d}u)^2+\dfrac{1}{K}\sin^2(\sqrt{K}u)(\mathrm{d}v)^2, & K>0,\\
(\mathrm{d}u)^2-\dfrac{1}{K}\sinh^2(\sqrt{-K}u)(\mathrm{d}v)^2, & K<0.
\end{cases}
\end{align*}
证明: 若 $(u(s),v(s))$ 是该曲面上的一条测地线的参数方程, 则存在不全为零的常数 $A,B,C$, 使得它按照 $K>0$ 还是 $K<0$ 分别满足下面的关系式:
\begin{align*}
&A\sin(\sqrt{K}u(s))\cos v(s)+B\sin(\sqrt{K}u(s))\sin v(s)+C\cos(\sqrt{K}u(s))=0,\\
&A\sinh(\sqrt{-K}u(s))\cos v(s)+B\sinh(\sqrt{-K}u(s))\sin v(s)+C\cosh(\sqrt{-K}u(s))=0.
\end{align*}
试把上述关系式和球面、伪球面上的测地线进行对照, 想一想: 上面的关系式有什么几何意义?
\end{example}

\begin{example}
在洛伦兹空间 $L^3$ 中求曲面 $\Sigma$ 在参数方程 (4.14) 下的自然标架的运动方程, 假定 $u^1=u,u^2=v$。设洛伦兹空间 $L^3$ 中经过原点的一个平面与曲面 $\Sigma$ 的交线的方程是 $\boldsymbol{r}=\boldsymbol{r}(s)$, 其中 $s$ 是弧长参数。证明:
\begin{align*}
\frac{\mathrm{d}^2\boldsymbol{r}(s)}{\mathrm{d}s^2}=-K\boldsymbol{r}(s).
\end{align*}
设曲线 $\boldsymbol{r}=\boldsymbol{r}(s)$ 用曲面 $\Sigma$ 的参数 $(u^1,u^2)$ 表示为 $u^1=u^1(s),u^2=u^2(s)$, 将上式改写为 $u^1(s),u^2(s)$ 所满足的常微分方程组。
\end{example}

\begin{example}
证明: 洛伦兹空间 $L^3$ 中经过原点的平面与伪球面 $\Sigma$ 的交线在球极投影 (4.13) 下的像 $(u(s),v(s))$ 满足方程
\begin{align*}
u^2(s)+v^2(s)-2au(s)-2bv(s)-\frac{4}{K}=0,
\end{align*}
其中 $a,b$ 是满足不等式 $\sqrt{a^2+b^2}>\frac{2}{\sqrt{-K}}$ 的常数。
\end{example}

\begin{example}
试求
Klein 圆: $u^2+v^2<1$, $\mathrm{d}s^2=\frac{(\mathrm{d}u)^2+(\mathrm{d}v)^2}{(1-(u^2+v^2))^2}$

和

Poincaré 上半平面: $y>0$, $\mathrm{d}s^2=\frac{1}{4y^2}(\mathrm{d}x^2+\mathrm{d}y^2)$

之间的保长对应。
\end{example}

\begin{example}
证明: 具有下列度量形式的曲面有常数 Gauss 曲率, 其值为 $-1$。试求它们之间的保长对应。
\begin{enumerate}[(1)]
\item $\mathrm{d}s^2=\frac{1}{v^2}\big((\mathrm{d}u)^2+(\mathrm{d}v)^2\big)$;
\item $\mathrm{d}s^2=(\mathrm{d}u)^2+\mathrm{e}^{2u}(\mathrm{d}v)^2$;
\item $\mathrm{d}s^2=(\mathrm{d}u)^2+\cosh^2 u(\mathrm{d}v)^2$.
\end{enumerate}
\end{example}



\end{document}